\lecture{23}{May 21.}{}
\begin{eg}[A dense proper subspace need not be closed]
    The space \(C^1 ([0, 1])\) of continuously differentiable functions on \([0, 1]\) is a linear subspace of \(C([0, 1])\). With respect to the supremum norm, \(C^1 ([0, 1])\) is dense in \(C([0, 1])\). For instance, every continuous function on \([0, 1]\) can be uniformly approximated by polynomials, and polynomials belong to \(C^1 ([0, 1])\). 

    However, \(C^1 ([0, 1]) \neq C([0, 1])\). For example, the function 
    \[
        f(t) \coloneqq  \left\vert t - \frac{1}{2} \right\vert 
    \]         
    is continuous on \([0, 1]\), but it is not differentiable at \(t = \frac{1}{2}\). Hence, \(C^1 ([0, 1])\) is a dense proper subspace of \(C ([0, 1])\). Now since we know a sequence of \(C^1 ([0, 1])\) does not converge in \(C^1 ([0, 1])\), \(C^1 ([0, 1])\) is not closed in 
    \[
        \left( C([0, 1]), \lVert \cdot \rVert_\infty \right). 
    \]     
    This shows that, in infitnite-dimensional normed spaces, linear subspaces need not be closed.
\end{eg}

\begin{eg}[Closed and bounded need not imply compact]
    Consider the Banach space 
    \[
        \left( C([0, 1]), \lVert \cdot \rVert_\infty  \right), 
    \]
    and let
    \[
        B \coloneqq \left\{ f \in C([0, 1]): \lVert f \rVert_\infty \le 1  \right\}. 
    \]
    Then \(B\) is closed and bounded. However, \(B\) is not compact. 

    To see this, define 
    \[
        f_n (t) \coloneqq \sin \left( 2^n \pi t \right), \quad 0 \le t \le 1. 
    \]  
    Then \(\lVert f_n \rVert_\infty \le 1 \), so \(f_n \in B\) for every \(n\). 

    We claim that the family \(\left\{ f_n: n \in \mathbb{N} \right\} \) is not equicontinuous. Indeed, take \(\varepsilon _0 = \frac{1}{2}\). Let \(\delta > 0\) be arbitrary. Since \(2^{-(n + 1)} \to 0\), we may choose \(n\) sufficiently large so that 
    \[
        \frac{1}{2^{n + 1}} < \delta.
    \]        
    Set 
    \[
        s_n = 0, \quad t_n = \frac{1}{2^{n + 1}}.
    \]
    Then 
    \[
        \vert s_n - t_n \vert = \frac{1}{2^{n + 1}} < \delta. 
    \]
    However, 
    \[
        f_n (s_n) = f_n (0) = 0,
    \]
    whereas 
    \[
        f_n (t_n) = \sin \left( 2^n \pi \cdot \frac{1}{2^{n+1} } \right) = \sin \left( \frac{\pi}{2} \right) = 1.  
    \]
    Therefore 
    \[
        \vert f_n (t_n) - f_n(s_n) \vert = 1 > \frac{1}{2} = \varepsilon_0. 
    \]
    Thus no single \(\delta > 0\) can make all the functions \(f_n\) uniformly continuous with the same modules of continuity. Hence \(\left\{ f_n \right\} \) is not equicontinuous. 

    Since \(\left\{ f_n: n \in \mathbb{N} \right\} \subseteq B\), it follows that \(B\) itself is not equicontinuous. By \autoref{cl: compactness criterion in C(K)}, \(B\) is not compact.        
\end{eg}

\begin{eg}[A linear functional need not be continuous]
    Let \((X, \lVert \cdot \rVert )\) be an infinite-dimensional normed vector space over \(\mathbb{R}\). We will construct a linear functional on \(X\) which is not continuous. 

    Choose a basis \(E\) of \(X\). By rescaling each basis vector if necessary, we may assume that \(\lVert e \rVert = 1\) for every \(e \in E\). This means that every non-zero vector \(x \in X\) can be written uniquely in the form 
    \[
        x = \sum_{i=1}^{\ell} x_i e_i, 
    \]       
    where \(e_1, \dots , e_{\ell } \in E\) are distinct basis vectors and \(x_i \in \mathbb{R} \setminus \left\{ 0 \right\} \). 

    Since \(X\) is infinite-dimensional, the basis \(E\) must be infinite.

    Because \(E\) is infinite, we may choose a sequence of distinct basis vectors \(e_1, e_2, e_3, \dots \in E\). Define a function \(\lambda : E \to \mathbb{R}\) by setting 
    \[
        \lambda (e_n) = n \quad \text{for } n = 1, 2, 3, \dots , 
    \]       
    and, for all other basis vectors \(e \in E \setminus \left\{ e_1, e_2, e_3, \dots  \right\} \), set \(\lambda (e) = 0\). Then \(\lambda \) is unbounded on \(E\). This means that there is no constant \(M \geq 0\) s.t.
    \[
        \vert \lambda (e) \vert \le M \quad \text{for every } e \in E.  
    \]     
    Indeed, if such an \(M\) existed, then for every positive integer \(n\) we would have 
    \[
        n = \vert \lambda (e_n) \vert \le M, 
    \]  
    which is impossible for all \(n\). Equivalently, for every \(M > 0\), one can choose \(n > M\) and then \(\vert \lambda (e_n) \vert = n > M \). Thus, the values of \(\lambda \) on the basis vectors become arbitrarily large. 

    We now use \(\lambda\) to define a linear functional on all of \(X\). If 
    \[
        x = \sum_{i=1}^{\ell } x_i e_i  
    \]       
    is the unique expansion of \(x\) with respect to the Hamel basis \(E\), define 
    \[
        \Phi_\lambda (x) \coloneqq \sum_{i=1}^{\ell } \lambda (e_i) x_i.  
    \]  
    This definition is well-defined because the representation of \(x\) as a finite linear combination of basis vectors is unique. 

    Moreover, \(\Phi _\lambda\) is linear. Indeed, the map is defined by prescribing its value on each basis vector,
    \[
        \Phi _\lambda (e) = \lambda (e), \quad e \in E,
    \]  
    and then extending linearly to finite linear combinations of basis vectors. Therefore, for all \(x, y \in X\) and all \(a, b \in \mathbb{R}\), 
    \[
        \Phi _\lambda (ax + by) = a \Phi _\lambda (x) + b \Phi _\lambda (y).
    \]  

    Now we claim that \(\Phi _\lambda\) is not bounded. Recall that a linear functional \(\Phi :X \to \mathbb{R} \) is bounded if there exists a constant \(C \ge 0\) s.t. 
    \[
        \vert \Phi (x) \vert \le C \lVert x \rVert \quad \text{for every } x \in X. 
    \]   
    In particular, if \(\Phi _\lambda\) were bounded, then for every basis vector \(e \in E\) we would have 
    \[
        \vert \lambda (e) \vert = \vert \Phi _\lambda (e) \vert \le C \lVert e \rVert = C
    \]  
    because \(\lVert e \rVert = 1\), which is impossible. 

    Finally, for linear maps between normed vector spaces, boundedness is equivalent to continuity. Hence \(\Phi _\lambda\) is a linear functional on \(X\) which is not continuous.  
\end{eg}

We now return to the compactness characterization of finite-dimensional normed spaces. Let \((X, \lVert \cdot \rVert )\) be a normed vector space, and denote its closed unit ball and unit sphere by 
\[
    B \coloneqq \left\{ x \in X: \lVert x \rVert \le 1  \right\}, \quad S \coloneqq \left\{ x \in X: \lVert x \rVert = 1  \right\}.  
\] 
In finite dimensions, the Heine-Borel theorem tells us \(B\) is compact. The result below shows that this property is not merely a consequence of finite-dimensionality, but actually characterizes it. 

\begin{lemma}[Riesz's lemma]
    Let \((X, \lVert \cdot \rVert )\) be a normed vector space, and let \(Y \subseteq X\) be a proper closed linear subspace. Then for every \(0 < \delta < 1\), there exists \(x \in X\) s.t. \(\lVert x \rVert = 1\) and
    \[
        \inf_{y \in Y} \lVert x - y \rVert \ge 1 - \delta.
    \]
    Equivalently, 
    \[
        \mathrm{dis}(x, Y) \ge 1 - \delta. 
    \]     
\end{lemma}

\begin{proof}
    Since \(Y \neq X\), we may choose a point \(x_0 \in X \setminus Y\). Geometrically, \(x_0\) is a point outside the closed subspace \(Y\). 

    Because \(Y\) is closed and \(x_0 \notin Y\), the distance from \(x_0\) to \(Y\) is strictly positive. Set 
    \[
        d \coloneqq \mathrm{dis}(x_0, Y) = \inf _{y \in Y} \lVert x_0 - y \rVert.  
    \]        
    Then \(d > 0\). The number \(d\) meausres how far \(x_0\) is from the subspace \(Y\). 

    Although there may not be a point of \(Y\) which realizes this distance, we can choose a point \(y_0 \in Y\) such that 
    \[
        \lVert x_0 - y_0 \rVert \le \frac{d}{1 - \delta},
    \]      
    which is because \(0 < \delta < 1\). 

    Now consider the vector from \(y_0 \in Y\) to \(x_0\): \(x_0 - y_0\). We normalize this vector and define 
    \[
        x \coloneqq \frac{x_0 - y_0}{\lVert x_0 - y_0 \rVert }. 
    \]    
    Then immediately \(\lVert x \rVert = 1\). Thus \(x\) is the unit vector pointing in the direction from the subspace \(Y\) toward the exterior point \(x_0\). 

    We claim that this unit vector remains at distance at least \(1 - \delta\) fro the whole subspace \(Y\). To prove this, let \(y \in Y\) be arbitrary. We want to estimate \(\lVert x - y \rVert \) from below. 

    Using the definition of \(x\), we have 
    \[
        x - y = \frac{x_0 - y_0}{\lVert x_0 - y_0 \rVert } - y.
    \]         
    Hence, 
    \begin{align*}
        \lVert x - y \rVert &= \frac{1}{\lVert x_0 - y_0 \rVert } \lVert x_0 - y_0 - \lVert x_0 - y_0 \rVert y  \rVert \\
        &= \frac{1}{\lVert x_0 - y_0 \rVert } \left\lVert x_0 - (y_0 + \lVert x_0 - y_0 \rVert y ) \right\rVert.  
    \end{align*}
    Now since \(y_0 + \lVert x_0 - y_0 \rVert y \in Y \), we have 
    \[
        \lVert x_0 - (y_0 + \lVert x_0 - y_0 \rVert y ) \rVert \ge \inf _{y \in Y} \lVert x_0 - y \rVert = d. 
    \] 
    Hence,
    \[
        \lVert x - y \rVert \ge \frac{d}{\lVert x_0 - y_0 \rVert }. 
    \]
    Now since 
    \[
        \lVert x_0 - y_0 \rVert \le \frac{d}{1 - \delta}, 
    \]
    we have 
    \[
        \lVert x - y \rVert \ge \frac{d}{\lVert x_0 - y_0 \rVert} \ge 1 - \delta. 
    \]
    Since \(y\) is arbitrarily chosen, we know 
    \[
        \inf _{y \in Y} \lVert x - y \rVert \ge 1 - \delta. 
    \] 
\end{proof}

We now apply Riesz's lemma to the compactness problem. If \(X\) is infinite-dimensional, then after choosing finitely many unit vectors \(x_1, \dots , x_k\), their span is still a proper finite-dimensional subspace of \(X\). Since finite-dimensional subspaces are closed, Riesz's lemma allows us to choose another unit vector \(x_{k + 1}\) that stays a fixed positive distance away from this span. Iterating this construction produces infinitely many unit vectors separated from one another by a uniform distance. Such a sequence cannot have a convergent subsequence, and therefore the unit sphere cannot be compact. 

This gives the desired converse to the finite-dimensional Heine-Borel theorem. 

\begin{theorem}[Compactness of the unit ball]
    Let \((X, \lVert \cdot \rVert )\) be a normed vector space. Then the following are equivalent:
    \begin{enumerate}[(i)]
        \item \(\dim X < \infty\). 
        \item The closed unit ball \(B\) is compact. 
        \item The unit sphere \(S\) is comapct.    
    \end{enumerate}
\end{theorem}

\begin{proof}
    We prove the implications \(\text{(i)} \implies \text{(ii)} \implies \text{(iii)} \implies \text{(i)}    \). 

    First assumte that \(\dim X < \infty\). By the finite-dimensional Heine-Borel theorem, a subset of \(X\) is compact if and only if it is closed and bounded. The closed unit ball 
    \[
        B = \left\{ x \in X: \lVert x \rVert \le 1  \right\} 
    \]  
    is closed and bounded. Therefore \(B\) is comapct. This proves \(\text{(i)} \implies \text{(ii)}  \). 

    Now assume that \(B\) is compact. The unit sphere 
    \[
        S = \left\{ x \in X: \lVert x \rVert = 1  \right\} 
    \]   
    is closed in \(X\). Indeed, the norm map 
    \[
        \lVert \cdot \rVert: X \to \mathbb{R}, \quad x \mapsto \lVert x \rVert 
    \] 
    is continuous, and \(S = \lVert \cdot \rVert^{-1} \left( \left\{ 1 \right\}  \right)  \). Since \(\left\{ 1 \right\} \) is closed in \(\mathbb{R}\), it follows that \(S\) is closed in \(X\). Since \(S \subseteq B\), the sphere \(S\) is a closed subset of the compact set \(B\). Hence, \(S\) is compact. Thus, (ii) \(\implies\) (iii).

    It remains to prove (iii) \(\implies\) (i). We prove the contrapositive. Suppose \(X\) is infinite-dimensional. We will show that \(S\) is not compact. 

    The idea is to construct a sequence of unit vectors that stay uniformly separated from one another. More precisely, we shall construct a sequence \((x_k)_{k=1}^{\infty} \subseteq S\) such that \(\lVert x_i - x_j \rVert \ge \frac{1}{2} \) whenever \(i \neq j\). Such a sequence cannot have a convergent subsequence, because every convergent sequence is Cauchy. Therefore, the existence of such a sequence implies that \(S\) is not compact. 

    Choose \(x_1 \in S\). Suppose that \(x_1, \dots , x_k \in S\) have already been chosen so that \(\lVert x_i - x_j \rVert \ge \frac{1}{2} \) for \(i \neq j\). Let \(Y_k \coloneqq \mathrm{span} \left\{ x_1, \dots , x_k \right\}  \). Then \(Y_k\) is finite-dimensional, hence closed in \(X\). Since \(X\) is infinite-dimensional, \(Y_k \neq X\). Thus \(Y_k\) is a proper closed linear subspace of \(X\). By Riesz's lemma, applied to \(Y_k\), there exists \(x_{k + 1} \in X\) such that \(\lVert x_{k+1} \rVert = 1 \) and \(\inf_{y \in Y_k} \lVert x_{k+1} - y \rVert \ge \frac{1}{2} \). Since \(x_1, \dots , x_k \in Y_k\), we have 
    \[
        \lVert x_{k+1} - x_i \rVert \ge \frac{1}{2} \quad \text{for } i = 1, \dots , k.  
    \]                                  
    Thus the induction continues. 
    
    Consequently, we obtain a sequence \((x_k) \subseteq S\) such that 
    \[
        \lVert x_i - x_j \rVert \ge \frac{1}{2} \quad \forall i \neq j, 
    \] 
    and we're done.
\end{proof}

\section{The Dual Space}
Let \(X\) and \(Y\) be normed vector spaces. Recall that \(\mathcal{L} (X, Y)\) denotes the vector space of all bounded linear operators \(A: X \to Y\), equipped with the operator norm 
\[
    \lVert A \rVert_{\mathcal{L} (X, Y)} \coloneqq \sup _{\lVert x \rVert_X \le 1} \lVert Ax \rVert_Y.  
\]
Equivalently,
\[
    \lVert Ax \rVert_Y \le \lVert A \rVert_{\mathcal{L} (X, Y)} \lVert x \rVert_X \quad \text{for all } x \in X.    
\]
The next theorem shows that \(\mathcal{L} (X, Y)\) is complete whenever the target space \(Y\) is complete. Notice that no completeness assumption is needed on \(X\). 

\begin{theorem}
    Let \(X\) be a normed vector space and let \(Y\) be a Banach space. Then \(\mathcal{L} (X, Y)\) is a Banach space with respect to the operator norm.   
\end{theorem}

\begin{proof}
    Let \((A_n)_{n \in \mathbb{N}}\) be a Cauchy sequence in \(\mathcal{L} (X, Y)\). Thus, for every \(\varepsilon > 0\), there exists \(n_0 \in \mathbb{N} \) s.t. 
    \[
        m, n \ge n_0 \implies \lVert A_m - A_n \rVert_{\mathcal{L} (X, Y)} < \varepsilon. 
    \]    
    For each fixed \(x \in X\), we have 
    \[
        \lVert A_n x - A_m x \rVert_Y = \lVert (A_n - A_m) x \rVert_Y \le \lVert A_n - A_m \rVert_{\mathcal{L} (X, Y)} \lVert x \rVert_X.    
    \] 
    Hence \((A_n x)\) is a Cauchy sequence in \(Y\). Since \(Y\) is complete, the limit 
    \[
        Ax \coloneqq \lim_{n \to \infty} A_n x 
    \]  
    exists for every \(x \in X\). This defines a map \(A: X \to Y\). 

    We first show that \(A\) is linear. Let \(x, x^{\prime} \in X\) and \(\alpha , \beta \in \mathbb{R}\). Then 
    \[
        A (\alpha x + \beta x^{\prime} ) = \lim_{n \to \infty} A_n (\alpha x + \beta x^{\prime} ) = \lim_{n \to \infty} \left( \alpha A_n x + \beta A_n x^{\prime}  \right) = \alpha \lim_{n \to \infty} A_n x + \beta \lim_{n \to \infty} A_n x^{\prime} = \alpha Ax + \beta Ay.     
    \]     
    Thus, \(A\) is linear. 

    It remains to show that \(A\) is bounded and that \(A_n \to A\) in operator norm. Fix \(\varepsilon > 0\). Since \((A_n)\) is Cauchy in operator norm, choose \(n_0 \in \mathbb{N}\) s.t. 
    \[
        m, n \ge n_0 \implies \lVert A_m - A_n \rVert_{\mathcal{L} (X, Y)} < \varepsilon. 
    \]      
    Fix \(n \ge n_0\). For every \(x \in X\), we have 
    \begin{align*}
        \lVert A x - A_n x \rVert_Y &= \lim_{m \to \infty} \lVert A_m x - A_n x \rVert_Y \le \limsup_{m \to \infty} \lVert A_m - A_n \rVert_{\mathcal{L} (X, Y)} \lVert x \rVert_X \le \varepsilon \lVert x \rVert_X.      
    \end{align*}  
    Therefore 
    \[
        \lVert (A - A_n) x \rVert_Y \le \varepsilon \lVert x \rVert_X \quad \text{for all } x \in X.   
    \]
    This implies \(A - A_n \in \mathcal{L} (X, Y)\) and 
    \[
        \lVert A - A_n \rVert_{\mathcal{L} (X, Y)} \le \varepsilon \quad \text{for every } n \ge n_0.  
    \] 
    In particular, taking \(n = n_0\), we get 
    \[
        \lVert Ax \rVert_Y \le \lVert Ax - A_{n_0} x \rVert_Y + \lVert A_{n_0} x \rVert_Y \le \left( \varepsilon + \lVert A_{n_0} \rVert_{\mathcal{L} (X, Y)}  \right) \lVert x \rVert_X.     
    \] 
    Hence \(A\) is bounded. 

    Moreover, 
    \[
        \lVert A - A_n \rVert_{\mathcal{L} (X, Y)} \le \varepsilon \quad \text{for all } n \ge n_0.  
    \]
    Thus \(A_n \to A\) in \(\mathcal{L} (X, Y)\).

    Since every Cauchy sequence in \(\mathcal{L} (X, Y)\) converges in \(\mathcal{L} (X, Y)\), \(\mathcal{L} (X, Y)\) is a Banach space.     
\end{proof}

A particularly important case is when the target space is \(\mathbb{R}\). The dual space of a normed vector space \(X\) is defined by 
\[
    X^* \coloneqq \mathcal{L} (X, \mathbb{R}).
\]  
Thus \(X^*\) is the space of all bounded linear functionals \(\Lambda : X \to \mathbb{R}\). By the preceding theorem, \(X^*\) is always a Banach space, whether or not \(X\) itself is complete. 

\begin{definition}[Dual space]
    Let \(X\) be a normed vector space. The dual space of \(X\) is \(X^* \coloneqq \mathcal{L} (X, \mathbb{R})\). Elements of \(X^*\) are called bounded linear functionals or continuous linear functionals on \(X\).     
\end{definition}

\begin{eg}[Dual space of \(L^p(\mu )\)]
    Let \((M, \mathcal{A} , \mu )\) be a measure space, and let \(1 < p < \infty\). Let \(q\) be the conjugate exponent of \(p\), that is, 
    \[
        \frac{1}{p} + \frac{1}{q} = 1.
    \]    
    Equivalently,
    \[
        q = \frac{p}{p - 1}, \quad q - 1 = \frac{1}{p - 1}, \quad (q - 1)p = q.
    \]
    Holder's inequality says that if \(f \in L^p (\mu)\) and \(g \in L^q (\mu)\), then \(f g \in L^1 (\mu)\) and 
    \begin{equation} \label{eq: 12.4}
        \left\vert \int _M fg \, \mathrm{d} \mu  \right\vert \le \int _M \vert f \vert \vert g \vert \, \mathrm{d} \mu \le \lVert f \rVert_p \lVert g \rVert_q. 
    \end{equation}  
    Consequently, every function \(g \in L^q (\mu)\) defines a linear functional 
    \[
        \Lambda _g : L^p (\mu) \to \mathbb{R}
    \] 
    by 
    \[
        \Lambda _g (f) \coloneqq \int _M f g \, \mathrm{d} \mu , \quad f \in L^p(\mu ). 
    \]
    The functional \(\Lambda _g\) is well-defined because \(fg \in L^1 (\mu)\). It is linear by the linearity of the integral. Moreover, by \autoref{eq: 12.4}, 
    \[
        \vert \Lambda _g (f) \vert \le \lVert f \rVert_p \lVert g \rVert_q \quad \text{for every } f \in L^p (\mu).
    \]   
    Therefore, \(\Lambda _g\) is bounded and 
    \begin{equation} \label{eq: 12.6}
        \lVert \Lambda _g \rVert_{\mathcal{L} \left( L^p(\mu ), \mathbb{R} \right) } \le \lVert g \rVert_q.
    \end{equation}
    In fact, the estimate is sharp. More precisely, one has 
    \[
       \lVert \Lambda _g \rVert_{\mathcal{L} \left( L^p(\mu ), \mathbb{R} \right) } = \lVert g \rVert_q. 
    \]
    We now prove the reverse inequality. If \(g = 0\), then \(\Lambda _g = 0\), and hence the equality is immediate. Thus assume \(g \neq 0\). Define 
    \[
        f \coloneqq \frac{\sgn (g) \vert g \vert^{q - 1} }{\lVert g \rVert_q^{q - 1} },
    \]   
    where, in the real-valued case, 
    \[
        \sgn (g) (x) = \begin{dcases}
            1, &\text{ if } g(x) > 0 ;\\
            -1, &\text{ if } g(x) < 0 ;\\
            0, &\text{ if } g(x) = 0.
        \end{dcases}
    \]
    Thus 
    \[
        \sgn (g) g = \vert g \vert, \quad \text{and} \quad \vert f \vert = \frac{\vert g \vert^{q - 1} }{\lVert g \rVert_q^{q - 1} }.    
    \]
    Since \((q - 1)p = q\), we have 
    \[
        \vert f \vert^p = \frac{\vert g \vert^{(q - 1) p} }{\lVert g \rVert_q^{(q - 1) p} } = \frac{\vert g \vert^q }{\lVert g \rVert_q^q }.
    \] 
    Because \(g \in L^q(\mu )\), this shows that \(f \in L^p(\mu )\). Moreover,
    \[
        \lVert f \rVert_p^p = \int _M \vert f \vert^p \, \mathrm{d} \mu = \frac{1}{\lVert g \rVert_q^q } \int _M \vert g \vert^q \, \mathrm{d} \mu = 1.     
    \]  
    Hence, \(\lVert f \rVert_p = 1 \). Now we evaluate \(\Lambda _g\) at this particular function \(f\). Using \(\sgn (g) g = \vert g \vert \), we get 
    \begin{align*}
        \Lambda _g (f) &= \int _M \frac{\sgn (g) \vert g \vert^{q - 1} }{\lVert g \rVert_q^{q - 1} } g \, \mathrm{d} \mu = \frac{1}{\lVert g \rVert_q^{q - 1} } \int _M \vert g \vert^q \, \mathrm{d} \mu = \frac{\lVert g \rVert_q^q }{\lVert g \rVert_q^{q - 1} } = \lVert g \rVert_q.    
    \end{align*} 
    Since \(\lVert f \rVert_p = 1 \), it follows from the definition of the operator norm that 
    \[
        \lVert \Lambda _g \rVert_{\mathcal{L} \left( L^p(\mu), \mathbb{R} \right) } \ge \left\vert \Lambda _g (f) \right\vert = \lVert g \rVert_q.   
    \]    
    Combining this lower bound with \autoref{eq: 12.6}, we obtain 
    \[
        \lVert \Lambda _g \rVert_{\mathcal{L} \left( L^p(\mu ), \mathbb{R}  \right) } = \lVert g \rVert_q.   
    \] 
\end{eg}

\section{Hilbert Space}
In the following, we introduce some elementary Hilbert space theory. The main result is the Riesz representation theorem, which shows that every Hilbert space is naturally isometrically isomorphic to its dual space. 

\begin{definition}[Innder Product] 
    Let \(H\) be a real vector space. A bilinear map \(H \times H \to \mathbb{R}\), \((x, y) \mapsto \langle x, y \rangle \), is called an inner product if it satisfies the following two properties: 
    \begin{enumerate}
        \item \textbf{Symmetry:} \(\langle x, y \rangle = \langle y, x \rangle  \) for all \(x, y \in H\). 
        \item \textbf{Positive definiteness:} \(\langle x, x \rangle \ge 0 \) for all \(x \in H\), and \(\langle x, x \rangle = 0 \) iff \(x = 0\).        
    \end{enumerate}   
\end{definition}

\begin{remark}
    The inner product induces a norm on \(H\) by 
    \[
        \lVert x \rVert \coloneqq \sqrt{\langle x, x \rangle }, \quad x \in H.  
    \] 
    This norm is called the norm associated with the inner product.
\end{remark}

\begin{lemma}[Cauchy-Schwarz inequality]
    Let \(H\) be a real vector space equipped with an inner product and the norm. Then the inner product and norm satisfy the Cauchy-Schwarz inequality 
    \[
        \vert \langle x, y \rangle  \vert \le \lVert x \rVert \lVert y \rVert   
    \] 
    and the triangle inequality 
    \[
        \lVert x + y \rVert \le \lVert x \rVert + \lVert y \rVert   
    \]
    for all \(x, y \in H\). Thus, the norm induce by the inner product is actually a norm. 
\end{lemma}

\begin{proof}
    The Cauchy-Schwarz inequality is obvious when \(x = 0\) or \(y = 0\). Hence assume that \(x, y \neq 0\). Define
    \[
        \xi \coloneqq \lVert x \rVert^{-1} x, \quad \eta \coloneqq  \lVert y \rVert^{-1} y.   
    \]   
    Then \(\lVert \xi \rVert = \lVert \eta  \rVert = 1  \). Hence 
    \[
        0 \le \lVert \eta - \langle \xi , \eta  \rangle \xi   \rVert^2 = \langle \eta , \eta  \rangle - \langle \xi, \eta \rangle^2 = 1 - \langle \xi , \eta  \rangle^2.  
    \] 
    This implies 
    \[
        \vert \langle \xi , \eta  \rangle \vert  \le 1. 
    \]
    Therefore 
    \[
        \vert \langle x, y \rangle  \vert \le \lVert x \rVert \lVert y \rVert.   
    \]

    It follows that 
    \begin{align*}
        \lVert x + y \rVert^2 = \lVert x \rVert^2 + 2 \langle x, y \rangle + \lVert y \rVert^2 \le \lVert x \rVert^2 + 2 \lVert x \rVert \lVert y \rVert + \lVert y \rVert^2 = \left( \lVert x \rVert + \lVert y \rVert   \right)^2.         
    \end{align*}
    Hence, we have the triangle equality 
    \[
        \lVert x + y \rVert \le \lVert x \rVert + \lVert y \rVert.   
    \]
\end{proof}

\begin{definition}[Hilbert space]
    Let \(H\) be a real vector space equipped with an inner product 
    \[
        \langle \cdot, \cdot \rangle : H \times H \to \mathbb{R}. 
    \] 
    The pair \((H, \langle \cdot, \cdot \rangle )\) is called a Hilbert space if \(H\) is complete with respect to the norm induced by the inner product, namely 
    \[
        \lVert x \rVert = \sqrt{\langle x, x \rangle }.  
    \]  
    Equivalently, \(H\) is a Hilbert space if every Cauchy sequence \(\left\{ x_n \right\}_{n=1}^{\infty} \) in \(H\), with respect to the metric 
    \[
        d(x, y) \coloneqq \lVert x - y \rVert, 
    \]  
    converges to an element of \(H\). 
\end{definition}

\begin{eg}[Dual space of a Hilbert space]
    Let \(H\) be a real Hilbert space. Thus \(H\) is a Banach space whose norm is induced by an inner product
    \[
        \langle \cdot, \cdot \rangle: H \times H \to \mathbb{R}, \quad \lVert x \rVert = \sqrt{\langle x, x \rangle }.   
    \]  
    Every element \(y \in H\) determines a linear functional \(\Lambda _y : H \to \mathbb{R}\) by 
    \[
        \Lambda _y (x) \coloneqq \langle x, y \rangle, \quad x \in H. 
    \]  
    This functional is bounded by the Cauchy-Schwarz inequality: 
    \[
        \vert \Lambda _y (x) \vert = \vert \langle x, y \rangle  \vert \le \lVert x \rVert \lVert y \rVert.    
    \]
    Hence, 
    \[
        \lVert \Lambda _y \rVert_{H^*} \le \lVert y \rVert. 
    \]
    In fact, taking \(x = y\) when \(y \neq 0\), or \(x = \frac{y}{\lVert y \rVert }\), shows that 
    \[
        \lVert \Lambda _y \rVert_{H^*} = \lVert y \rVert.  
    \]
    The Riesz representation theorem states that every bounded linear functional on \(H\) is of this form. Therefore the map \(H \to H^*\), \(y \mapsto \Lambda _y\), is an isometric isomorphism.   
\end{eg}

\begin{eg}
    Let \((M, \mathcal{A}, \mu )\) be a measure space. Then \(H \coloneqq L^2 (\mu)\) is a Hilbert space. The inner product is induced by the bilinear map 
    \[
        \mathcal{L} ^2 (\mu ) \times \mathcal{L} ^2 (\mu ) \to \mathbb{R}, \quad (f, g) \mapsto \langle f, g \rangle \coloneqq \int _M fg \, \mathrm{d} \mu. 
    \]  
    It is well-defined because the product of two \(L^2\)-functions \(f, g: M \to \mathbb{R}\) is integrable by the Cauchy-Schwarz inequality. That it is bilinear and symmetric follows directly from the properties of Lebesgue integral.  
\end{eg}

\begin{theorem}[Riesz representation theorem]
    Let \(H\) be a Hilbert space and let \(\Lambda : H \to \mathbb{R}\) be a bounded linear functional. Then there exists a unique element \(y \in H\) such that 
    \[
        \Lambda (x) = \langle y, x \rangle \quad \text{for all } x \in H.  
    \]  
    This element \(y \in H\) satisfies 
    \[
        \lVert y \rVert = \sup _{0 \neq x \in H} \frac{\vert \langle y, x \rangle  \vert }{\lVert x \rVert } = \lVert \Lambda \rVert.  
    \] 
    Thus the map 
    \[
        H \to H^*, \quad y \mapsto \langle y, \cdot \rangle 
    \]
    is an isometry of normed vector spaces.
\end{theorem}

\begin{remark}
    Riesz's representation theorem shows \(H \simeq H^*\) for any Hilbert space \(H\).  
\end{remark}

\begin{definition}[Convex set]
    Let \(V\) be a real vector space. A subset \(C \subseteq V\) is called convex if for every \(x, y \in C\) and every \(t \in [0, 1]\), one has 
    \[
        tx + (1 - t)y \in C.
    \]    
    Equivalently, \(C\) contains the line segment joining any two of its points. 
\end{definition}

\begin{theorem}[Projection theorem] \label{thm: projection theorem}
    Let \(H\) be a Hilbert space and let \(E \subseteq H\) be a nonempty closed convex subset. Then there exists a unique elements \(x_0 \in E\) s.t. 
    \[
        \lVert x_0 \rVert \le \lVert x \rVert \quad \text{for all } x \in E.   
    \]   
\end{theorem}

\begin{proof}
    Define 
    \[
        \delta \coloneqq \inf \left\{ \lVert x \rVert: x \in E  \right\}. 
    \]
    We first prove uniqueness. Suppose that \(x_0, x_1 \in E\) both satisfy \(\lVert x_0 \rVert = \lVert x_1 \rVert = \delta  \). Since \(E\) is convex, 
    \[
        \frac{x_0 + x_1}{2} \in E.
    \] 
    Therefore,
    \[
        \left\lVert \frac{x_0 + x_1}{2} \right\rVert \ge \delta. 
    \]
    Using the parallelogram identity, we get 
    \[
        \lVert x_0 - x_1 \rVert^2 = 2 \lVert x_0 \rVert^2 + 2 \lVert x_1 \rVert^2 - \lVert x_0 + x_1 \rVert^2 \le 4 \delta ^2 - 4 \delta ^2 = 0.    
    \]
    Hence, \(x_0 = x_1\). 

    Now we prove the existence. Choose a sequence \((x_i) \subseteq E\) s.t. \(\lim_{i \to \infty} \lVert x_i \rVert = \delta  \). We claim that \((x_i)\) is a Cauchy sequence. 

    Fix \(\varepsilon > 0\), \(\exists i_0 \in \mathbb{N}\) s.t. 
    \[
        i \ge i_0 \implies \lVert x_i \rVert^2 < \delta ^2 + \frac{\varepsilon^2}{4}. 
    \]      
    Let \(i, j \in \mathbb{N}\) be such that \(i \ge i_0, j \ge i_0\). Since \(E\) convex, \(\frac{x_i + x_j}{2} \in E\). Hence,
    \[
        \left\lVert \frac{x_i + x_j}{2} \right\rVert \ge \delta. 
    \]    
    Using parallelogram identity again,
    \[
        \lVert x_i - x_j \rVert^2 = 2 \lVert x_i \rVert^2 + 2 \lVert x_j \rVert^2 - \lVert x_i + x_j \rVert^2 < 4 \left( \delta ^2 + \frac{\varepsilon^2}{4} \right) - 4 \delta ^2 = \varepsilon ^2.     
    \]
    Thus, \(\lVert x_i - x_j \rVert < \varepsilon\) for all \(i, j \ge i_0\)m and thus \((x_i)\) is a Cauchy sequence. 

    Now since \(H\) is complete, the limit \(x_0 \coloneqq \lim_{i \to \infty} x_i \) exists in \(H\). Moreover, \(x_0 \in E\) because \(E\) is closed. Finally, by continuity of the norm,  \(\lVert x_0 \rVert = \delta \). Hence,
    \[
        \lVert x_0 \rVert \le \lVert x \rVert \quad \text{for all } x \in E.   
    \]         
\end{proof}